\chapter{Observability, Evaluation and Testing}
\label{ch:observability}

\begin{quote}\itshape
Without an evaluation set you are not changing a prompt. You are guessing and
shipping.
\end{quote}

\noindent
Three subjects, one argument: you cannot improve what you cannot see, and you
cannot change what you cannot re-measure. This is the chapter that separates
people who have run an agent in production from people who have built one.

% ============================================================
\section{Tracing}
\label{sec:tracing}
\index{observability!tracing}

Tracing switches on with environment variables --- \code{LANGSMITH\_TRACING},
an API key, a project name --- and from that point every \api{Runnable} and every
graph node appears in a trace without further code. That low activation cost is
the whole reason to do it before you think you need it.

What a trace gives you that logs do not: the prompt \emph{as sent} after every
middleware has modified it, the tool arguments the model actually chose, token
counts per step, and latency attributed per node rather than per request.

\code{@traceable} adds your own functions to the same trace, which matters
because the interesting question is usually ``what did our retrieval do'' rather
than ``what did the model do''.

\begin{note}
\textbf{When a client will not accept a SaaS backend} --- and in regulated
sectors many will not --- OpenTelemetry, Langfuse and Phoenix are the
alternatives, the first being the one that keeps you portable. You lose the
tighter integration with datasets and evaluation runs; you keep the traces, which
is the part you cannot reconstruct later.

Do not let this become a reason to have no tracing. Chapter~\ref{ch:asyncio}'s
blocking call and Chapter~\ref{ch:multi-agent}'s handoff loop are both nearly
undiagnosable without one.
\end{note}

\subsection{The metrics worth a dashboard}
\label{sec:metrics}

\begin{center}
\small
\begin{tabularx}{\linewidth}{@{}lX@{}}
\toprule
\textbf{Metric} & \textbf{Why} \\
\midrule
Cost per conversation & The number that decides whether the product works.
  Chapter~\ref{ch:context} moved it 5$\times$. \\[3pt]
Time to first token & What a user experiences as speed
  (Chapter~\ref{ch:interrupts}). \\[3pt]
Steps per run & Rises before anything breaks. The earliest warning of a
  loop. \\[3pt]
Tool error rate, per tool & One flaky tool degrades everything downstream of
  it. \\[3pt]
Task completion & The only one that measures the product rather than the
  system. \\
\bottomrule
\end{tabularx}
\end{center}

% ============================================================
\section{Evaluation}
\label{sec:evaluation}
\index{evaluation}

The surface lives in \pkg{langsmith}:

\begin{python}[caption={An evaluation run}, label={lst:evaluate}]
from langsmith import evaluate

results = evaluate(
    target,                      # what to run
    data="ops-copilot-v1",       # the dataset
    evaluators=[called_the_right_tool, schema_valid, judged_relevance],
    num_repetitions=5,           # run each case N times
    max_concurrency=8,
)
\end{python}

Also present: \code{aevaluate} for the async twin, \code{evaluate\_existing} to
re-score runs you already have, and \code{evaluate\_comparative} for A/B.

\begin{note}
\code{num\_repetitions} is a first-class parameter, and its presence is a
statement. Running each case once tells you what happened once. Running it five
times tells you the \emph{distribution}, which for a non-deterministic system is
the only honest summary --- and it is exactly the discipline
Chapter~\ref{ch:multi-agent} insisted on for the topology benchmark. The API
agrees with the book.
\end{note}

\mermaidfig{eval-ladder}{%
  Four tiers by cost and confidence. Push work downwards.}{%
  The evaluation ladder from deterministic checks to human review.}

\subsection{Trajectory, not just the answer}
\label{sec:trajectory}

This is the agent-specific part and the part most teams skip. Two agents can
produce the same final text having done completely different things --- one
called the right tool, one guessed; one took three steps, one took eleven and
cost four times as much.

\begin{center}
\small
\begin{tabularx}{\linewidth}{@{}lX@{}}
\toprule
\textbf{Trajectory check} & \textbf{Catches} \\
\midrule
Was the expected tool called? & Silent regression when a description changes \\[3pt]
Were unexpected tools called? & Scope creep, and injection
  (Ch.~\ref{ch:security}) \\[3pt]
Step count within bounds? & Loops, before the recursion limit fires \\[3pt]
Cost within budget? & The regression nobody notices until the invoice \\
\bottomrule
\end{tabularx}
\end{center}

All four are deterministic, free, and computable from a trace. They belong in CI.

\subsection{Judges, and why an uncalibrated one is worse than none}
\label{sec:judges}
\index{evaluation!LLM as judge}

For qualities you cannot check mechanically --- relevance, groundedness, tone ---
a model scores the output. It works, and it has three failure modes worth naming:
judges prefer longer answers, prefer their own family's style, and drift when the
judge model is upgraded underneath you.

\begin{warning}
\textbf{Calibrate against human labels and record the agreement rate.} Hand-label
thirty cases, run the judge, and compute how often it agrees. If you cannot state
that number, your evaluation set produces numbers of unknown meaning --- and
because they look precise, they will be trusted more than a subjective opinion
that at least announces itself as one.

Pin the judge model explicitly, and re-calibrate when you change it. A judge
upgraded silently moves every historical score.
\end{warning}

\subsection{In CI}
\label{sec:eval-ci}

The mechanism is easy: run the set, compare against a threshold, fail the build.
The difficulty is flakiness --- a genuinely non-deterministic system will cross a
tight threshold occasionally for no reason, and a gate that cries wolf gets
disabled within a month.

Three things that make it survivable: use \code{num\_repetitions} and threshold
on the median rather than a single run; set the threshold as a delta from the
current baseline rather than an absolute; and separate ``blocks the merge'' from
``posts a comment''. Only the deterministic checks should block.

% ============================================================
\section{Testing}
\label{sec:testing}
\index{testing!fake models}

\mermaidfig{test-pyramid-agents}{%
  What runs on every commit, and what is deliberately excluded.}{%
  The testing layers for an agent codebase.}

The claim that non-determinism makes agents untestable is wrong, and precisely
wrong: it makes \emph{text assertions} untestable. Assert on behaviour --- which
tool was called, whether the schema validates, how many steps it took --- and the
suite is as deterministic as any other.

\begin{python}[caption={The seam everything hangs on}, label={lst:fake}]
class ScriptedFakeChatModel(FakeMessagesListChatModel):
    """No in-box fake implements bind_tools; see Chapter 10."""

    def bind_tools(self, tools, **kwargs):
        return self
\end{python}

Scripted to emit specific tool calls in a specific order, that fake tests graph
topology exactly --- which branch was taken, what the reducer merged, whether the
approval gate fired --- with no provider and no variance.

\begin{center}
\small
\begin{tabularx}{\linewidth}{@{}lX@{}}
\toprule
\textbf{Layer} & \textbf{Tested how} \\
\midrule
Tools & Plain functions. Ordinary unit tests. \\[3pt]
Reducers & Pure two-argument functions --- the cheapest test in the codebase
  (Ch.~\ref{ch:stategraph}) \\[3pt]
Nodes & \code{state -> dict}. No graph needed. \\[3pt]
Graph flow & Scripted fake model \\[3pt]
Quality & Evaluation set, not unit tests \\[3pt]
Provider integration & A handful, marked, excluded from the default run \\
\bottomrule
\end{tabularx}
\end{center}

\begin{warning}
Chapter~\ref{ch:cancellation}'s trap applies with force here: if
\code{asyncio\_mode} is lost, every \code{async def} test reports \textbf{passed}
without running. A suite can be green and testing nothing. Have a test whose only
job is to detect that.
\end{warning}

% ============================================================
\section{Prompts as configuration}
\label{sec:prompt-versioning}

Changing a prompt without a deploy is attractive and is a trap in exactly one
respect: it decouples the change from the thing that would have caught it. A
prompt in a database that anyone can edit is a production change with no review,
no rollback and no evaluation run.

If you do it --- and there are good reasons to --- version the prompts, record
which version served each run in the trace, and make the evaluation set part of
the promotion path rather than something that runs afterwards.

% ============================================================
\section{What to take away}
\label{sec:ch14-summary}

\begin{itemize}[leftmargin=1.4em]
\item Tracing costs environment variables. Turn it on before you think you need
      it; the failures it diagnoses are otherwise nearly invisible.
\item \code{num\_repetitions} is first-class because one run of a
      non-deterministic system tells you nothing about the distribution.
\item Trajectory checks --- right tool, step count, cost --- are deterministic,
      free, and the layer most teams skip.
\item An uncalibrated judge produces precise-looking numbers of unknown meaning.
      Record the agreement rate; pin the judge model.
\item Non-determinism does not make agents untestable. It makes text assertions
      untestable.
\item Only deterministic checks should block a merge.
\item A green suite that makes no assertions is the failure mode to test for
      explicitly.
\end{itemize}

\begin{exercisebox}
\textbf{1.} Build a twenty-case set with three evaluator kinds. Run it with
\code{num\_repetitions=1}, then \code{=5}. Compare the variance and decide which
you would gate on.

\textbf{2.} Hand-label thirty outputs, run a judge over the same thirty, and
compute the agreement. That number is the credibility of every judged score you
will report.

\textbf{3.} Deliberately worsen a prompt and confirm the gate catches it. An
untested gate is not a gate.

\textbf{4.} Write a trajectory evaluator that asserts the expected tool was
called and the step count stayed under six. Add it to CI --- it needs no
provider.
\end{exercisebox}

\begin{projectbox}
\textbf{Stage 11 --- observability, evaluation, testing.} Tracing wired up, a
twenty-case evaluation set with a regression threshold in CI, and the full suite
consolidated.

The contract's non-negotiable line: \textbf{the default suite makes no real model
calls}. If that has stopped being true, the stage has failed regardless of what
else passes --- because the book spends a chapter arguing for it and the project
exists to demonstrate that the argument is practical.
\end{projectbox}
